\section{Introduction}\label{sec:intro}

Scientific writing is not a neutral container for ideas. The way results are framed, in particular, how directly arguments are stated, how much jargon is used, whether evidence is foregrounded, and how strongly novelty is claimed, can shape both peer evaluation and the dissemination of knowledge. This possibility echoes a broader concern, associated with the idea of a ``different voice" \parencite{gilligan1993different}, that dominant evaluative standards may overlook systematically different modes of communication. Yet empirical work in economics has focused much more on gender differences in publication, promotion, and review outcomes than on the writing itself.

This paper asks whether the gender composition of author teams is associated with systematic differences in how economists write. The paper builds most directly on \textcite{hengel2022publishing}, who shows that female-authored papers are more readable and appear to face higher writing standards in peer review. We extend that insight beyond readability by measuring a broader set of surface-level writing features in abstracts, including directness, jargon, evidence and citation usage, active voice, novelty claims, hedging, and related rhetorical dimensions. The objective is not to infer author intent or personality from text. Rather, it is to identify whether different teams systematically present research in different ways.

Our data come from \textcite{hengel2022publishing}, with the main analysis conducted on 9,117 abstracts from four top general-interest economics journals between 1950 and 2015. Each abstract is evaluated using a 16-item LLM-based rubric covering readability, technicality, directness, evidentiary support, and rhetorical style. Because LLM-based measures may invite skepticism, the empirical analysis proceeds in several steps. First, we benchmark the LLM readability score against the readability measures used in \textcite{hengel2022publishing}. Second, we examine unconditional mean differences across all 16 criteria. Third, we estimate regressions that add journal, year, article, and author controls to assess whether the observed patterns persist after accounting for observable compositional differences. Finally, we extend the analysis to peer review, testing both whether peer review differentially affects writing by author gender and whether any such differential erodes as female authors accumulate publishing experience.

The validation exercise shows that the LLM readability measure tracks the same underlying pattern documented by \textcite{hengel2022publishing}. In unconditional comparisons, female-majority teams score higher on readability, directness, accessibility of language, and evidence/citation usage, and they make weaker novelty claims. Once controls are added, the clearest differences remain higher readability, less jargon, greater directness, more evidentiary support, and weaker novelty claims among female-majority teams. By contrast, differences in assertiveness, qualifiers, emotional valence, and several other tonal dimensions are small or unstable across specifications. The results, therefore, do not support a simple story in which women write in a uniformly more hesitant or emotional way. The clearest differences are concentrated instead in clarity, technical density, and the organization of claims.

We then extend this analysis to how the differences evolve during peer review and with cumulative publishing experience. Here, we show that while peer review measurably affects LLM-evaluated writing quality on average, it does so similarly for male- and female-authored papers, in contrast to the differential effect \textcite{hengel2022publishing} document for readability. In tests of cumulative exposure, the aggregate ``preemptive adaptation'' pattern Hengel identifies (in which experienced female authors increasingly produce cleaner drafts rather than relying on review to close the gap) is largely absent across the LLM scores. Patterns resembling that mechanism do, however, resurface in three specific dimensions. Acknowledgement of limitations and modal verb strength both exhibit a pattern resembling the preemptive-adaptation signature, which would be consistent with female authors learning and internalizing reviewer expectations over time. Novelty-claim strength instead shows a distinct pattern in which peer review consistently dampens female authors' novelty framing across all experience levels; one reading of this is persistent reviewer resistance rather than rational adaptation. We emphasize that these dimension-specific readings are suggestive, not definitive: evaluation noise and the limited granularity of the 1 to 10 rubric remain plausible contributors to the observed patterns.

This paper contributes to three literatures. First, it adds to work on gender and evaluation in economics by studying writing style as a potential input into evaluation rather than focusing only on publication or career outcomes. Second, it extends readability analysis from \textcite{hengel2022publishing} to a multidimensional measurement framework that separates clarity and structure from tone and affect. Third, by showing which dimensions move together with team gender composition and which do not, and where the peer-review mechanism Hengel documents is and is not visible, it sharpens the debate over what a ``different voice" in economics actually consists of and where differential review pressure may bind most tightly.

The remainder of the paper proceeds as follows. Section~\ref{sec:lit} reviews the related literature. Section~\ref{sec:data_and_measu} describes the data, authorship measures, and LLM rubric. Section~\ref{sec:method} outlines the empirical strategy. Section~\ref{sec:results} presents the validation exercise, mean-difference analysis, and controlled regressions. Section~\ref{sec:conclusion} concludes.